These findings are subject to several limitations. The benchmark was restricted to two simulated-driving datasets, so generalisation to other experimental paradigms or clinical populations was not tested. All analyses were also performed offline on recorded data, so performance under real driving conditions, where movement and electrode contact may differ, was not assessed.
Performance was not reported separately by blink duration, so detection of prolonged eye closures, which is particularly relevant to drowsiness assessment, remains unestablished. The evaluation depends on manually annotated reference labels and on their definitions of blink onset and offset, which can vary with signal representation and annotation criteria~\citep{nystrm2024what}. \rev{The Internal labels were anchored to synchronised video and the eye-aspect-ratio signal, whereas the Cao2018 labels relied on EEG morphology alone; this difference in annotation modality limits the comparability of absolute scores across the two corpora.} In addition, the second rater reviewed the first rater's annotations rather than annotating independently from scratch, so inter-rater agreement was not quantified.

%\rev{Several additional methodological limitations warrant disclosure. A matched factorial comparison varying only the calibration pool while holding other pipeline components constant was not performed, so the improvement cannot be attributed exclusively to epoch screening. In the present implementation, the recording is epoched before per-epoch filtering and resampling, which differs from the continuous-first order described in Section~II and may introduce duration-dependent edge effects. The blink-enrichment of candidate epochs selected by Stage~A was not quantified, so the extent to which the selected pool is specific to blink-related activity, as opposed to other high-amplitude artefacts, remains unestablished. The median-versus-mean Stage~B comparison (Proposed-Mean versus Proposed-approach) yielded nonsignificant differences on both corpora, consistent with the centre estimator not being the primary driver of the performance advantage. The headline strategy comparison used a retrospective best-channel-per-session rule based on reference $F_1$; this common oracle was applied to every condition, but the resulting scores do not represent prospective channel selection without labels. Finally, sessions from the same participant were treated as independent experimental units, which may understate the uncertainty of the paired comparisons; participant-clustered or mixed-effects analyses are needed to assess the sensitivity of the statistical conclusions to this simplification.}

